% ============================================================================
% PHẦN 3 (2023 - 2024): TỐI ƯU HÓA CHÍNH SÁCH TRỰC TIẾP (DPO & CÁC BIẾN THỂ)
% ============================================================================
\sectionintro{3}{DPO và các biến thể}{2023--2024}{closed-form solution, gradient, kiến trúc hai mô hình và các hướng cải tiến.}

\subsection{3.1. Từ RLHF Objective đến DPO Loss (closed-form solution)}

% ----------------------------------------------------------------------------
% SLIDE: ĐỘNG CƠ CỦA DPO (05/2023)
% ----------------------------------------------------------------------------
\begin{frame}{DPO: tối ưu trực tiếp từ sở thích}
  \lead{Thay vòng lặp RL bốn mô hình bằng closed-form solution (Rafailov et al., 2023).}
  \subhead{Liên hệ trực tiếp Reward Function $r(x,y)$ và Policy $\pi^*(y\mid x)$}
  \small
  \begin{itemize}\setlength{\itemsep}{4pt}
    \item \textbf{Quy trình gián tiếp của RLHF + PPO (Phần 2):}
      \[
        \fitmath{%
        \text{Cặp } (y_w, y_l) \xrightarrow{\text{Bước 1 (MLE)}} \text{Reward Function } r_\phi(x, y) \xrightarrow{\text{Bước 2 (PPO + } V_\psi\text{)}} \text{Policy } \pi_\theta(y \mid x)
        }
      \]
    \item \textbf{Quy trình trực tiếp của DPO (Phần 3):} Dùng closed-form solution của bài toán RLHF để biểu diễn $r(x, y)$ theo $\pi_\theta(y \mid x)$, từ đó tối ưu trực tiếp $\pi_\theta$ trên cặp mẫu dương/âm $(y_w, y_l)$!
  \end{itemize}

  \vspace{0.06cm}
  \small
  \renewcommand{\arraystretch}{1.18}
  \begin{tabularx}{\textwidth}{@{}l X X@{}}
    \toprule
    \textbf{Tiêu chí Hệ thống} & \textbf{RLHF + PPO (Phần 2 -- 2022)} & \textbf{DPO (Phần 3 -- 2023)} \\
    \midrule
    \textbf{Số mô hình trong VRAM} & \textbf{4 mạng:} $\pi_\theta, \pi_{\text{ref}}, r_\phi, V_\psi$ & \textbf{2 mạng:} Chỉ cần $\pi_\theta$ và $\pi_{\text{ref}}$ \\

    \textbf{Autoregressive Decoding lúc học} & Bắt buộc ở mọi bước cập nhật (\textit{Rollout}) & \textbf{Không cần} (Tính log-prob song song) \\
    \bottomrule
  \end{tabularx}
  \par\vspace{0pt}
\end{frame}

% ----------------------------------------------------------------------------
% SLIDE: TỪ MỤC TIÊU RLHF (B1) ĐẾN DPO Loss (B7)
% ----------------------------------------------------------------------------
\begin{frame}{DPO: closed-form solution cho RLHF}
  \lead{Bài toán có KL regularization và phân phối Gibbs--Boltzmann (B1).}
\subhead{RLHF Objective và closed-form solution (B1)}
      \small
      \textbf{Bài toán cực trị có KL regularization:}
      \[
        \fitmath{%
        \max_{\pi} \E_{y \sim \pi} \left[ r(x, y) - \beta \log \frac{\pi(y \mid x)}{\pi_{\text{ref}}(y \mid x)} \right]
        }
      \]
      \textbf{Nghiệm cực trị Gibbs--Boltzmann (BĐT Gibbs):}
      \[
        \fitmath{%
        \pi^*(y \mid x) = \frac{1}{Z(x)} \pi_{\text{ref}}(y \mid x) \exp\left(\frac{1}{\beta} r(x, y)\right)
        }
      \]
      \textbf{Biểu diễn ngược Reward Function:}
      \[
        \fitmath{%
        \boxed{r(x, y) = \beta \log \frac{\pi^*(y \mid x)}{\pi_{\text{ref}}(y \mid x)} + \beta \log Z(x)}
        }
      \]
      \textit{\footnotesize($Z(x) = \sum_{y} \pi_{\text{ref}}\exp(\frac{1}{\beta}r)$: Partition Function)}
  \par\vspace{0pt}
\end{frame}

\begin{frame}{DPO: khử Partition Function và học từ cặp mẫu}
  \lead{Từ closed-form solution của RLHF đến Loss trực tiếp (B7).}
\subhead{Khử $Z(x)$ và xây dựng DPO (B7)}
      \small
      \textbf{Thế vào xác suất cặp Bradley--Terry:}\\
      Khi lấy hiệu số $r(x, y_w) - r(x, y_l)$, đại lượng $\beta \log Z(x)$ \textbf{tự triệt tiêu hoàn toàn}:
      \[
        \fitmath{%
        r(x, y_w) - r(x, y_l) = \beta \log \frac{\pi^*(y_w \mid x)}{\pi_{\text{ref}}(y_w \mid x)} - \beta \log \frac{\pi^*(y_l \mid x)}{\pi_{\text{ref}}(y_l \mid x)}
        }
      \]
      \vspace{-0.15cm}
      \begin{block}{DPO Loss (B7)}
        \footnotesize
        Thay $\pi^*$ bằng $\pi_\theta$, tối ưu MLE trên tập cặp:
        \[
          \fitmath{%
          \mathcal{L}_{\text{DPO}}(\theta) = -\E_{(x, y_w, y_l)} \left[ \log \sigma \left( \hat{r}_\theta(x, y_w) - \hat{r}_\theta(x, y_l) \right) \right]
          }
        \]
        với phần thưởng ngầm định:
        $\hat{r}_\theta(x, y) = \beta \log \frac{\pi_\theta(y \mid x)}{\pi_{\text{ref}}(y \mid x)}$.
      \end{block}
  \par\vspace{0pt}
\end{frame}

\subsection{3.2. Cơ chế Gradient, Kiến trúc 2 Mô hình \& Các Biến thể (2024)}

% ----------------------------------------------------------------------------
% SLIDE: CƠ CHẾ GRADIENT VÀ TỰ ĐIỀU TRỌNG SỐ CỦA DPO
% ----------------------------------------------------------------------------
\begin{frame}{Gradient của DPO}
  \lead{Tự điều chỉnh biên độ cập nhật giữa mẫu dương $y_w$ và mẫu âm $y_l$.}
  \subhead{Gradient đầy đủ (Rafailov et al., 2023)}
  \small
  \[
    \fitmath{%
    \nabla_\theta \mathcal{L}_{\text{DPO}}(\theta) = -\beta \E_{(x, y_w, y_l)} \Big[ \underbrace{\sigma\big(\hat{r}_\theta(x, y_l) - \hat{r}_\theta(x, y_w)\big)}_{\text{Hệ số lỗi xếp hạng } \in (0, 1)} \cdot \Big( \underbrace{\nabla_\theta \log \pi_\theta(y_w \mid x)}_{\text{Kéo Mẫu dương } y_w \; (\uparrow)} - \underbrace{\nabla_\theta \log \pi_\theta(y_l \mid x)}_{\text{Đẩy Mẫu âm } y_l \; (\downarrow)} \Big) \Big]
    }
  \]

  \subhead{Trọng số thay đổi theo lỗi xếp hạng}
  \footnotesize
  \begin{itemize}\setlength{\itemsep}{4pt}
    \item \textbf{Khi mô hình xếp hạng sai ($\hat{r}_\theta(x, y_l) > \hat{r}_\theta(x, y_w)$):}\\
      Hệ số lỗi $\sigma(\hat{r}_l - \hat{r}_w) > 0.5$ (tiệm cận $1.0$), tự động \textbf{tăng biên độ gradient} để sửa lỗi xếp hạng.
    \item \textbf{Khi mô hình đã phân loại đúng ($\hat{r}_\theta(x, y_w) \gg \hat{r}_\theta(x, y_l)$):}\\
      Hệ số $\sigma(\hat{r}_l - \hat{r}_w) \to 0$, gradient tự triệt tiêu $\to$ \textbf{ngăn ngừa quá khớp} trên các mẫu đã học tốt.
    \item \textbf{Tác động đối ngẫu:}\\
      Đồng thời tăng hợp lý của mẫu dương $y_w$ và giảm xác suất của mẫu âm $y_l$.
  \end{itemize}
  \par\vspace{0pt}
\end{frame}

% ----------------------------------------------------------------------------
% SLIDE: SƠ ĐỒ KIẾN TRÚC & LUỒNG DỮ LIỆU CỦA DPO
% ----------------------------------------------------------------------------
\begin{frame}{DPO: kiến trúc hai mô hình}
  \lead{Chỉ giữ Policy $\pi_\theta$ và Reference Model $\pi_{\mathrm{ref}}$.}
  \vspace{0.15cm}
  \begin{center}
    \resizebox{0.96\textwidth}{!}{%
      \begin{tikzpicture}
        \node[data, text width=3.3cm, minimum height=1.3cm] (data) at (0, 0) {
          \textbf{Cặp đối chiếu tĩnh}\\
          $(x, y_w, y_l) \sim \mathcal{D}_{\text{pair}}$\\
          \normalfont\tiny ($y_w$: Mẫu dương, $y_l$: Mẫu âm)
        };

        \node[trainmodel, text width=3.8cm, minimum height=1.35cm] (actor) at (5.0, 1.2) {
          \textbf{1. Policy Model $\pi_\theta$}\\
          (Đang tối ưu hóa)\\
          \normalfont\scriptsize Tính $\log \pi_\theta(y_w \mid x), \;\log \pi_\theta(y_l \mid x)$
        };

        \node[frozenmodel, text width=3.8cm, minimum height=1.35cm] (ref) at (5.0, -1.2) {
          \textbf{2. Reference Model $\pi_{\text{ref}}$}\\
          (Cố định tham số)\\
          \normalfont\scriptsize Tính $\log \pi_{\text{ref}}(y_w \mid x), \;\log \pi_{\text{ref}}(y_l \mid x)$
        };

        \node[rewardbox, text width=3.7cm, minimum height=2.1cm] (implicit) at (10.1, 0) {
          \textbf{Thưởng Ngầm định}\\[4pt]
          $\hat{r}_w = \beta \log \dfrac{\pi_\theta(y_w \mid x)}{\pi_{\text{ref}}(y_w \mid x)}$\\[4pt]
          $\hat{r}_l = \beta \log \dfrac{\pi_\theta(y_l \mid x)}{\pi_{\text{ref}}(y_l \mid x)}$
        };

        \node[lossbox, text width=3.5cm, minimum height=2.1cm] (dpoloss) at (14.8, 0) {
          \textbf{DPO Loss}\\[4pt]
          $\mathcal{L}_{\text{DPO}} = -\log \sigma(\hat{r}_w - \hat{r}_l)$\\[4pt]
          \normalfont\scriptsize Khử hoàn toàn $r_\phi$\\
          và Value Network $V_\psi$
        };

        \coordinate (bus_in) at (2.25, 0);
        \draw[line width=1.05pt, draw=navyblue!85!black] (data.east) -- (bus_in);
        \draw[flow] (bus_in) |- (actor.west);
        \draw[flow] (bus_in) |- (ref.west);

        \draw[flow] (actor.east) -- (7.55, 1.2) |- ([yshift=0.45cm]implicit.west);
        \draw[flow] (ref.east) -- (7.55, -1.2) |- ([yshift=-0.45cm]implicit.west);
        \draw[flow] (implicit.east) -- (dpoloss.west);

        \draw[gradflow] (dpoloss.north) -- (14.8, 2.45) -- node[above, font=\scriptsize\rmfamily\bfseries, text=stanfordred] {Backpropagation trực tiếp cập nhật $\pi_\theta$ (không cần Rollout)} (5.0, 2.45) -- (actor.north);
      \end{tikzpicture}%
    }
  \end{center}
  \vspace{0.05cm}
  \small
  \textbf{Ưu thế hệ thống:} Tính toán log-xác suất song song trên toàn chuỗi $(y_w, y_l)$ trong một forward pass duy nhất.
  \par\vspace{0pt}
\end{frame}

% ----------------------------------------------------------------------------
% SLIDE: CÁC BIẾN THỂ HẬU DPO (2024: SIMPO, RPO, IPO, KTO)
% ----------------------------------------------------------------------------
\begin{frame}{Các giới hạn của DPO chuẩn}
  \lead{Length Bias, Likelihood Displacement và yêu cầu dữ liệu cặp.}
\limithead{Ba hạn chế của DPO}
      \small
      \begin{enumerate}\setlength{\itemsep}{4pt}
        \item \textbf{Length Bias:}\\
          Tổng $\log \pi_\theta(y \mid x) = \sum_{t=1}^{|y|} \log \pi_\theta(y_t \mid x, y_{<t})$ tỷ lệ thuận với $|y|$, gây ưu tiên mẫu dài dòng.
        \item \textbf{Likelihood Displacement:}\\
          Chỉ tối ưu hiệu số $\hat{r}_w - \hat{r}_l$, có thể làm cả $\log \pi_\theta(y_w)$ lẫn $\log \pi_\theta(y_l)$ cùng giảm.
        \item \textbf{Đòi hỏi cặp đồng truy vấn $(x, y_w, y_l)$:}\\
          Không tận dụng được dữ liệu đánh giá nhị phân đơn lẻ.
      \end{enumerate}
  \par\vspace{0pt}
\end{frame}

\begin{frame}{Các biến thể của DPO (2024)}
  \lead{SimPO, RPO / CPO, IPO và KTO xử lý những giới hạn khác nhau.}
\motivehead{Các hướng cải tiến (2024)}
      \small
      \begin{itemize}\setlength{\itemsep}{4pt}
        \item \textbf{SimPO (Meng et al., 2024):}\\
          Chuẩn hóa độ dài $\frac{\beta}{|y|}\log \pi_\theta(y \mid x)$ và thay $\pi_{\text{ref}}$ bằng hằng số biên $\gamma > 0$ (chỉ cần \textbf{1 mô hình}).
        \item \textbf{RPO / CPO (2024):}\\
          Cộng thêm $\alpha \mathcal{L}_{\text{SFT}}(y_w)$ để neo xác suất tuyệt đối của mẫu dương $y_w$.
        \item \textbf{IPO (Azar et al., 2024):}\\
          Loss bình phương ngăn quá khớp khi xác suất cặp tiệm cận $1$.
        \item \textbf{KTO (Ethayarajh et al., 2024):}\\
          Học trực tiếp từ mẫu dương hoặc mẫu âm đơn lẻ.
      \end{itemize}
  \par\vspace{0pt}
\end{frame}

\subsection{3.3. Tổng kết Phần 3 (DPO): Thành tựu \& Các Vấn đề Tồn đọng cho Phần 4}

% ----------------------------------------------------------------------------
% SLIDE: TỔNG KẾT PHẦN 3 (1/2): NHỮNG GÌ ĐÃ ĐẠT ĐƯỢC & TỐI ƯU SO VỚI RLHF+PPO
% ----------------------------------------------------------------------------
\begin{frame}{DPO thay đổi hệ thống RLHF như thế nào?}
  \lead{Đối chiếu DPO và các biến thể (2023--2024) với PPO (2022).}
  \motivehead{Đóng góp so với RLHF + PPO}
  \vspace{0.05cm}
  \small
  \renewcommand{\arraystretch}{1.20}
  \begin{tabularx}{\textwidth}{@{}p{3.4cm} X X@{}}
    \toprule
    \textbf{Tiêu chí đánh giá} & \textbf{RLHF + PPO (Phần 2 -- 2022)} & \textbf{DPO \& Biến thể (Phần 3 -- 2023--2024)} \\
    \midrule
    \textbf{1. Cơ sở giải tích} & Xấp xỉ lặp qua gradient Monte Carlo và mạng $r_\phi, V_\psi$ & \textbf{\textcolor{tealgreen}{Boltzmann closed-form solution}}: Khử hoàn toàn Partition Function $\beta \log Z(x)$ \\

    \textbf{2. Bộ nhớ GPU} & \textbf{4 mạng:} $\pi_\theta, \pi_{\text{ref}}, r_\phi, V_\psi$ & \textbf{\textcolor{tealgreen}{2 mạng}} ($\pi_\theta, \pi_{\text{ref}}$) với DPO, hoặc \textbf{1 mạng} với SimPO \\

    \textbf{3. Tốc độ tính toán} & Chậm do Autoregressive Decoding tuần tự (\textit{Rollout}) & \textbf{\textcolor{tealgreen}{Tăng tốc vượt bậc}}: Truyền thuận song song trên cặp $(y_w, y_l)$ \\

    \textbf{4. Tính ổn định} & Phương sai gradient cao, nhạy cảm siêu tham số & Hội tụ ổn định tương đương SFT (ứng dụng trên Llama 3, Qwen 2, Zephyr) \\
    \bottomrule
  \end{tabularx}
  \par\vspace{0pt}
\end{frame}

% ----------------------------------------------------------------------------
% SLIDE: TỔNG KẾT PHẦN 3 (2/2): VẤN ĐỀ TỒN ĐỌNG MÀ PHẦN 4 (RLVR + GRPO) GIẢI QUYẾT
% ----------------------------------------------------------------------------
\begin{frame}{Giới hạn của DPO trong bài toán suy luận}
  \limithead{Vì sao quay lại Online RL vào cuối 2024?}
  \small
  \begin{enumerate}\setlength{\itemsep}{2pt}
    \item \textbf{Off-policy Exploration Bottleneck:}\\
      \small DPO chỉ tối ưu trên tập cặp tĩnh $\mathcal{D}_{\text{pair}}$ thu thập trước; không thể tự khám phá các quỹ đạo suy luận mới nằm ngoài tập dữ liệu.
    \item \textbf{Định giá bằng mạng nơ-ron trên toán học và thuật toán:}\\
      \small Trong bài toán suy luận đa bước, cả DPO lẫn mạng $r_\phi$ đều bị đánh lừa bởi các lời giải trình bày đẹp nhưng sai kết quả (\textit{Reward Hacking}).
    \item \textbf{Điểm nghẽn Value Network $V_\psi$ khi Quay lại Online RL:}\\
      \small Nếu dùng lại PPO (Phần 2) cho chuỗi suy luận dài $10{,}000$ token, Critic $V_\psi$ bị suy biến do phần thưởng thưa ở cuối chuỗi.
  \end{enumerate}

  \vspace{0.06cm}
  \begin{exampleblock}{Tiếp theo: RLVR và GRPO (cuối 2024--2026)}
    \small
    Kết hợp \textbf{Deterministic Verifier (RLVR)} chống \textit{Reward Hacking} và giải thuật \textbf{GRPO} chuẩn hóa nhóm $G$ quỹ đạo để \textbf{xóa bỏ Critic $V_\psi$}.
  \end{exampleblock}
  \par\vspace{0pt}
\end{frame}
